%======================================================================
\chapter{Stairwells, and Theorem 20 from Theorem 15}\label{ch:stairwells}
%======================================================================

This chapter reduces Theorem~\ref{thm:ho20} to a statement about separators alone (Hoehn--Oversteegen's Theorem~15, proved in
Chapter~\ref{ch:ladders}). A separator $M$ of the bottom from the top of $G^\star$ is approximated from inside a neighbourhood $U$ by a
\emph{stairwell}: finitely many sheets $S_1,\dots,S_k$, each the graph of a function over a subgraph of $G$, joined end to end, with $k$ odd. Think of a
staircase that climbs from left to right over $G$, turns back, climbs over a smaller part, turns again, and so on. If $k=1$ the single sheet is a section
of $G^\star\to G$, and $f$ lifts into it exactly. If $k>1$, one simple fold of $G$ unwinds one turn of the staircase: on the fold, the preimage of the
stairwell contains a stairwell whose ``pit'' has moved up one level. Theorem~\ref{thm:ho8} lifts $f$ to the fold, and after finitely many folds the height
is $1$.

Throughout, $G$ is a PL graph in a normed space $E$, $\pi\colon E\times\R\to E$ the projection, $G^\star=G\times\I$; relative frontiers $\rfr$ are taken in $G$. For a simple fold
$F$ on $G$ with projection $\varphi$ we write $F^\star=F\times\I\subseteq(E\times\R)\times\R$, $\varphi^\star(q,t)=(\varphi(q),t)$, and $\varphi^{\star-1}(S)$ for the preimage of $S\subseteq G^\star$ in
$F^\star$. Hoehn and Oversteegen's ``regular subsets'' of $G$ are read as PL subgraphs throughout.

\section{Sides}

\begin{definition}[Hoehn--Oversteegen, Definition~9]\label{def:sides}
Let $A\subseteq G$ be a PL graph and $B\subseteq\rfr A$.
\begin{enumerate}[label=(\alph*)]
\item $A$ has \emph{consistent complement relative to $B$} if for every component $C$ of $G\setminus A$ either $\rfr C\subseteq B$ or $\rfr C\cap B=\emptyset$.
\item The \emph{$A$ side of $B$}, $\side_B(A)$, is the closure of the union of the components of $G\setminus B$ that meet $A$.
\end{enumerate}
\end{definition}

\formal{\lean{HPC.HO16.HasConsCompl}, \lean{HPC.HO16.sideOf}.}

Everything below uses that $G\setminus K$ has finitely many components for a PL graph $K\subseteq G$ or a finite $K$, each open in $G$ with PL closure and
frontier inside $\rfr K$ (Proposition~\ref{prop:pl-library}).

\formal{\lean{HPC.HO16.IsPLGraph.finite_components_diff}, \lean{HPC.HO16.IsPLGraph.comp_diff_mem_nhdsWithin},
\lean{HPC.HO16.IsPLGraph.closure_comp_diff}, \lean{HPC.HO16.IsPLGraph.relFrontier_comp_diff}.}

\begin{proposition}[the observations after Definition~9]\label{prop:sides-char}
Let $G$ be connected, $A\subseteq G$ a PL graph and $B\subseteq\rfr A$. In (b)--(d) assume that $A$ and $B$ are nonempty and that $A$ has consistent complement relative to $B$.
\begin{enumerate}[label=(\alph*)]
\item If $B=\emptyset$ then $\side_B(A)$ is the union of the components of $G$ meeting $A$; in particular $\side_B(\emptyset)=\emptyset$.
\item For every neighbourhood $V$ of $B$, $\side_B(A)=\side_B(A\cap V)$.
\item $\side_B(A)$ is a PL graph containing $A$, $\rfr\side_B(A)=B$, and $\side_B(A)\cap V=A\cap V$ for some neighbourhood $V$ of $B$.
\item $\side_B(A)$ is the only closed $D\subseteq G$ with $\rfr D=B$ and $D\cap V=A\cap V$ for some neighbourhood $V$ of $B$.
\end{enumerate}
\end{proposition}

\begin{proof}
(b) Let $K$ be a component of $G\setminus B$ meeting $A$. As $G$ is connected and $B\neq\emptyset$, $\cl K$ meets $B$ at some $b$. A component $C$ of $G\setminus A$ inside $K$
has $\rfr C\cap B=\emptyset$: otherwise $\rfr C\subseteq B$, so $C$ is clopen in $K$ and $C=K$, but $K$ meets $A$. There are finitely many such $C$, none with $b$ in its closure;
so the points of $K$ near $b$ are in $A$, and $K$ meets $A\cap V$.

(c) $\rfr\side_B(A)\subseteq B$ since the union of components is open. Conversely $b\in B$ is a limit of points of $A\setminus B$ ($B$ is finite and $A$ has no isolated
points), so $b\in\side_B(A)$; and $b$ is in the closure of a component $C$ of $G\setminus A$, which then has $\rfr C\subseteq B$, so $C$ is a component of $G\setminus B$
missing $A$, and $C\cap\side_B(A)=\emptyset$. For $V$ take a neighbourhood of $B$ missing the closures of the finitely many components $C$ of $G\setminus A$ with
$\rfr C\cap B=\emptyset$.

(d) If $\rfr D=B$, then $D\setminus B$ is clopen in $G\setminus B$, hence a union of components; each of them meets $V\cap V'$ near a point of its closure in $B$, so it
meets $\side_B(A)\setminus B$; and conversely.
\end{proof}

\formal{\lean{HPC.HO16.ho16_obs9_emptyB}, \lean{HPC.HO16.ho16_obs9_nhds}, \lean{HPC.HO16.ho16_obs9_char}, \lean{HPC.HO16.ho16_obs9_unique}; the
common core is \lean{HPC.HO16.IsPLGraph.sideOf_char_aux}.}

\begin{proposition}[Hoehn--Oversteegen, Proposition~10]\label{prop:ho10}
Let $G$ be connected, $A,A'\subseteq G$ nonempty PL graphs, $B\subseteq\rfr A\cap\rfr A'$, both of consistent complement relative to $B$, and $A\cap V=A'\cap V$ for a
neighbourhood $V$ of $B$. Then $\side_B(A)=\side_B(A')$; and every component $C$ of $G\setminus A$ with $\cl C\cap B\neq\emptyset$ is a component of $G\setminus A'$.
\end{proposition}

\begin{proof}
The first part is Proposition~\ref{prop:sides-char}(d). For the second, $\rfr C\subseteq B$ by consistency. Shrink $V$ so that it misses the closures of the finitely
many components of $G\setminus A'$ whose closures miss $B$. Let $C'$ be a component of $G\setminus A'$ meeting $C\cap V$. Then $C'\subseteq C$ (since $\rfr C\subseteq B$ and $C'\cap
B=\emptyset$), and $\rfr C'$ meets $B$, so $\rfr C'\subseteq B$. If $C'\neq C$, a point of $\rfr C'\cap C$ would lie in $B\cap C\subseteq A\cap C=\emptyset$.
\end{proof}

\formal{\lean{HPC.HO16.ho16_prop10}.}

In Hoehn and Oversteegen's proof of the second part, consistency of $A'$ is applied to a component $C'$ that merely meets $C\cap V$; that its frontier meets
$B$ needs the shrinking of $V$ above. The second part is used again in Proposition~\ref{prop:ho19}.

\begin{proposition}[Hoehn--Oversteegen, Proposition~11]\label{prop:ho11}
Let $G$ be connected, $A\subseteq G$ a nonempty PL graph, and $\rfr A=B_1\cup B_2$ with $B_1\cap B_2=\emptyset$, $A$ of consistent complement relative to $B_1$ and $B_2$.
Then $(\side_{B_1}(A),A,\side_{B_2}(A))$ is a fold triple on $G$, with $\rfr\side_{B_1}(A)=B_1$ and $\rfr\side_{B_2}(A)=B_2$. Consistency relative to $B_1$ implies
consistency relative to $B_2$.
\end{proposition}

\begin{proof}[Proof sketch]
Every component of $G\setminus A$ has frontier inside $B_1$ or inside $B_2$; $\side_{B_1}(A)$ is $A$ together with the closures of the components whose frontier lies in
$B_2$, and symmetrically. This gives $\side_{B_1}(A)\cup\side_{B_2}(A)=G$ and $\side_{B_1}(A)\cap\side_{B_2}(A)=A$; the closures of the two differences are separated
because their attaching points lie in the disjoint sets $B_2$ and $B_1$. The frontiers come from Proposition~\ref{prop:sides-char}(c).
\end{proof}

\formal{\lean{HPC.HO16.ho16_prop11}, \lean{HPC.HO16.ho16_prop11_frontier}, \lean{HPC.HO16.ho16_prop11_remark}.}

\section{Straight sets and stairwell structures}

\begin{definition}[Hoehn--Oversteegen, Definition~12]\label{def:straight}
\begin{enumerate}[label=(\alph*)]
\item A point $x$ of $G$ is \emph{ordinary} if some neighbourhood of $x$ in $G$ is the union of two segments $[a,x]$, $[x,b]$ meeting only in $x$. (So $x$ is neither an end point
nor a branch point of $G$.)
\item A list $\langle B_1,\dots,B_n\rangle$ of finite subsets of $G$ is \emph{generic} if all their points are ordinary and the $B_i$ are pairwise disjoint.
\item $S\subseteq G^\star$ is \emph{straight} if $S$ is closed, $\pi$ is injective on $S$, and $\pi(S)$ is a PL graph. Its \emph{end set} is $\mathcal E(S)=S\cap\pi^{-1}(\rfr\pi(S))$.
\end{enumerate}
\end{definition}

\formal{\lean{HPC.HO16.IsOrdinaryPt}, \lean{HPC.HO16.IsGeneric}, \lean{HPC.HO16.IsStraight}, \lean{HPC.HO16.endSet}.}

A straight set need not be a PL graph: it is the graph of a continuous function over the PL graph $\pi(S)$. Ordinariness does not refer to a decomposition:
for any decomposition, $x$ is ordinary iff it is interior to an edge or the end of exactly two edges; all but finitely many points are ordinary. $\pi$ maps
$S\setminus\mathcal E(S)$ openly into $G$.

\formal{\lean{HPC.HO16.isOrdinaryPt_iff_order}, \lean{HPC.HO16.IsPLGraph.finite_not_ordinary}, \lean{HPC.HO16.ho16_obs12}.}

\begin{definition}[Hoehn--Oversteegen, Definition~13]\label{def:stairwell}
A \emph{stairwell structure of height $k\ge1$} for $S\subseteq G^\star$ is a list $\langle S_1,\dots,S_k\rangle$ with finite sets $\alpha_1,\dots,\alpha_{k+1}\subseteq G^\star$ (write
$\beta_i=\alpha_{i+1}$) such that:
\begin{enumerate}[label=(S\arabic*)]
\item $S_1,\dots,S_k$ are nonempty straight sets with union $S$;
\item $\mathcal E(S_i)=\alpha_i\cup\beta_i$ with $\alpha_i\cap\beta_i=\emptyset$, and $\alpha_1=\emptyset=\beta_k$;
\item for $i<k$ there is a neighbourhood $V$ of $\pi(\beta_i)$ with $\pi(S_i)\cap V=\pi(S_{i+1})\cap V$;
\item each $\pi(S_i)$ has consistent complement relative to $\pi(\alpha_i)$ and to $\pi(\beta_i)$;
\item $\langle\pi(\alpha_2),\dots,\pi(\alpha_k)\rangle$ is generic in $G$.
\end{enumerate}
\end{definition}

\formal{\lean{HPC.HO16.IsStairwellStr} (with the sets $\alpha_i$ as one sequence, so that $\beta_i=\alpha_{i+1}$ holds by construction) and its existential form
\lean{HPC.HO16.HasStairwell}.}

So consecutive sheets share their end points and agree near them over $G$; sheet $i$ meets sheet $i+1$ only at $\beta_i$. The sets $\alpha_i$ may be empty for
$2\le i\le k$.

\begin{proposition}[remarks on pp.~193--194]\label{prop:stair-remarks}
Let $\langle S_1,\dots,S_k\rangle$ be a stairwell structure for $S$.
\begin{enumerate}[label=(\alph*)]
\item If $C$ is a component of $G$ meeting every $\pi(S_i)$, intersecting all sets with $C^\star$ gives a stairwell structure for $S\cap C^\star$ relative to the graph $C$.
\item If $G$ is connected and $k=1$, then $\pi$ maps $S$ homeomorphically onto $G$.
\item \emph{(Claim 14.1.)} If $G$ is connected and $k$ odd, every $x\in G$ lies in an odd number of the sets $\pi(S_i)$; in particular $\pi(S)=G$.
\item Without connectedness, (b) and (c) fail: for $G=[0,1]\cup[2,3]$, $S=[0,1]\times\set{\frac12}$ is a stairwell structure of height $1$ (its end set is empty, as $\rfr[0,1]=\emptyset$ in
$G$) and $\pi(S)\neq G$.
\end{enumerate}
\end{proposition}

\begin{proof}
(a) $C$ is clopen in $G$, so frontiers, components of complements and ordinary points relative to $C$ are those relative to $G$. (b) $\mathcal E(S_1)=\emptyset$, so
$\pi(S_1)$ is a nonempty clopen subset of $G$, so it is $G$; and $\pi|_{S_1}$ is a continuous injection of a compactum. (c) By induction along the stairwell, using
Proposition~\ref{prop:ho10} with (S3) and Proposition~\ref{prop:ho11} with $B_1=\pi(\alpha_i)$, $B_2=\pi(\beta_i)$: crossing $\pi(\beta_i)$ changes the count by $0$ or $2$.
\end{proof}

\formal{\lean{HPC.HO16.ho16_obs13_comp}, \lean{HPC.HO16.ho16_height_one}, \lean{HPC.HO16.ho16_claim14_1}, \lean{HPC.HO16.ho16_odd_image},
\lean{HPC.HO16.ho16_p193_remark_false}.}

Hoehn and Oversteegen allow disconnected graphs (``Throughout the remainder of this paper, $G$ will denote a (not necessarily connected) graph''), and state the
remarks (b), (c) on pp.~193--194 without connectedness; (d) shows that it is needed there. Their Proposition~14, to which the remarks refer, does assume it.

\begin{definition}[Hoehn--Oversteegen, Definition~16]\label{def:broken}
A \emph{broken stairwell structure of height $k$ with a pit at level $i_0$} ($1\le i_0\le k$) for $S\subseteq G^\star$ is a list $\langle S_1,\dots,S_k;P_1,P_2\rangle$, finite sets
$\alpha_1,\dots,\alpha_{k+1}$ (with $\beta_i=\alpha_{i+1}$) and a finite set $\gamma$ such that:
\begin{enumerate}[label=(S\arabic*$'$)]
\item $S_1,\dots,S_k,P_1,P_2$ are nonempty straight sets with union $S$;
\item (S2) holds except that $\mathcal E(S_{i_0})=\alpha_{i_0}\cup\beta_{i_0}\cup\gamma$, a disjoint union; and $\mathcal E(P_2)=\mathcal E(P_1)\cup\gamma$ with $\mathcal E(P_1)\cap\gamma=\emptyset$;
\item (S3) holds, and $\pi(P_1)$ and $\pi(P_2)$ agree near $\pi(\mathcal E(P_1))$, $\pi(P_2)$ and $\pi(S_{i_0})$ agree near $\pi(\gamma)$;
\item (S4) holds, $\pi(S_{i_0})$ has consistent complement relative to $\pi(\gamma)$, and $\pi(P_2)$ relative to $\pi(\mathcal E(P_1))$ and to $\pi(\gamma)$;
\item $\langle\pi(\alpha_2),\dots,\pi(\alpha_k),\pi(\mathcal E(P_1)),\pi(\gamma)\rangle$ is generic;
\item $\pi(\alpha_{i_0})\cap\pi(P_1\cup P_2)=\emptyset$.
\end{enumerate}
\end{definition}

\formal{\lean{HPC.HO16.IsBrokenStairwellStr}; the restriction to a component is \lean{HPC.HO16.ho16_obs16_comp}.}

The pit is the pair $P_1,P_2$: a short staircase hanging from sheet $i_0$ at $\gamma$, going back over $\pi(P_2)$ and forward again over $\pi(P_1)$, ending at $\mathcal E(P_1)$.

\begin{proposition}[Hoehn--Oversteegen, Proposition~17]\label{prop:ho17}
If $S$ has a stairwell structure $\langle S_1,\dots,S_k\rangle$ of height $k\ge3$, then $\langle S_3,\dots,S_k;S_1,S_2\rangle$ is a broken stairwell structure for $S$ of height $k-2$ with a pit at level~$1$,
with $\alpha_1'=\emptyset$, $\alpha_i'=\alpha_{i+2}$ for $i\ge2$, and $\gamma'=\alpha_3$.
\end{proposition}

\begin{proof}
$\mathcal E(P_2)=\mathcal E(S_2)=\alpha_2\cup\alpha_3=\mathcal E(P_1)\cup\gamma'$, since $\alpha_1=\emptyset$ gives $\mathcal E(P_1)=\alpha_2$; and $\alpha_2\cap\alpha_3=\emptyset$. The neighbourhoods of (S3$'$) are those of
(S3) at $i=1,2$; (S4$'$) is (S4) at $i=2,3$; the list of (S5$'$) is a permutation of that of (S5); and (S6$'$) holds as $\alpha_1'=\emptyset$.
\end{proof}

\formal{\lean{HPC.HO16.ho16_prop17}, \lean{HPC.HO16.ho16_prop17_exists}. The hypothesis $k\ge3$ is not printed; for $k\le2$ the conclusion has height $\le0$.}

\begin{proposition}[remarks on p.~201]\label{prop:rem17}
Let $\langle S_1,\dots,S_k;P_1,P_2\rangle$ be a broken stairwell structure for $S$. If $\gamma=\emptyset$, then $\langle S_1,\dots,S_k\rangle$ is a stairwell structure for $S_1\cup\dots\cup S_k$. If
$\mathcal E(P_1)=\emptyset$ and $G$ is connected, then $\langle P_1\rangle$ is a stairwell structure of height~$1$ for $P_1$.
\end{proposition}

\formal{\lean{HPC.HO16.ho16_rem17_gamma}, \lean{HPC.HO16.ho16_rem17_P1}.}

\section{Pulling straight sets back to a fold}

\begin{lemma}[Hoehn--Oversteegen, Lemma~18]\label{lem:ho18}
Let $F=F_1\cup F_2\cup F_3$ be a simple fold on $G$, and $S\subseteq G^\star$ straight. Suppose that $\rfr\pi(S)\cap\rfr\varphi(F_2)=\emptyset$, or that $\varphi(F_2)\cap V\subseteq\pi(S)\cap V$ for a
neighbourhood $V$ of $\rfr\pi(S)\cap\rfr\varphi(F_2)$. Then, with relative frontiers of subsets of $F$ taken in $F$:
\begin{enumerate}[label=(\arabic*)]
\item $S'=\varphi^{\star-1}(S)$ is straight in $F^\star$, and $\mathcal E(S')=\varphi^{\star-1}(\mathcal E(S))\setminus(\rfr F_2)^\star$;
\item $S''=S'\cap(F_1\cup F_2)^\star$ is straight in $F^\star$, and $\mathcal E(S'')=\bigl([\varphi^{\star-1}(\mathcal E(S))\cap(F_1\cup F_2)^\star]\setminus(\rfr F_1)^\star\bigr)\cup(S''\cap(\rfr F_3)^\star)$.
\end{enumerate}
\end{lemma}

\begin{proof}
\emph{Claim: if $x\in S'$ lies over $\rfr F_2$, then $\varphi(W)\subseteq\pi(S)$ for a neighbourhood $W$ of $\pi(x)$ in $F$.} Say $\pi(x)\in\rfr F_1=F_1\cap F_2$, so $y=\varphi(\pi(x))\in\rfr G_1\subseteq\rfr G_2$
and $y\in\pi(S)$. If $y$ is interior to $\pi(S)$ take $V\subseteq\pi(S)$; otherwise the hypothesis gives $V$. Shrink $V$ to miss $\cl{G_1\setminus G_2}$, so that $G_1\cap V=G_2\cap V$.
$F_3$ is closed and misses $\pi(x)$, so $W=\varphi^{-1}(V)\cap(F_1\cup F_2)$ is a neighbourhood of $\pi(x)$ in $F$, and $\varphi(W)=G_2\cap V\subseteq\pi(S)$.

(1) $\pi(S')=\varphi^{-1}(\pi(S))$, and $\pi$ is injective on $S'$ because it is on $S$. Off $(\rfr F_2)^\star$, $\varphi$ is a local homeomorphism (Lemma~\ref{lem:ho6}(7)), so
frontier points of $\pi(S')$ correspond to frontier points of $\pi(S)$; over $\rfr F_2$ the claim puts $\pi(x)$ in the interior of $\pi(S')$. $\pi(S')$ is a PL graph plus finitely many
points, and has no isolated points, so it is a PL graph.

(2) The same on $F_1\cup F_2$; in addition a point over $\rfr F_3=F_2\cap F_3$ is an end point, since every neighbourhood of it meets $F_3\setminus F_2$, which misses $\pi(S'')$.
\end{proof}

\formal{\lean{HPC.HO16.ho16_lemma18}, with the preimage \lean{HPC.HO16.foldPre} and $\varphi^\star$ as \lean{HPC.HO16.phiStar}; PL-graph facts for preimages:
\lean{HPC.HO16.IsPLGraph.inter_preimage_fst}, \lean{HPC.HO16.isPLGraph_union_finite}.}

Hoehn and Oversteegen check that the components of $\pi(S')$ are nondegenerate but not that there are finitely many. On p.~202 they print
``$\mathcal E(S'')\subset(F_1\cap F_2)^\star$ since $S''\subset(F_1\cap F_2)^\star$'' for $(F_1\cup F_2)^\star$.

\section{Proposition 19: unwinding one turn}

\begin{proposition}[Hoehn--Oversteegen, Proposition~19]\label{prop:ho19}
Let $G$ be connected and $S\subseteq G^\star$ have a broken stairwell structure $\langle S_1,\dots,S_k;P_1,P_2\rangle$ with a pit at level $i_0$, in which $\gamma\neq\emptyset$ and $\mathcal E(P_1)\neq\emptyset$.
Then there is a simple fold $F$ on $G$ with $F$ connected such that $\varphi^{\star-1}(S)$ contains a set $S'$ with a broken stairwell structure of height $k$ with a pit at level $i_0+1$ if
$i_0<k$, and a stairwell structure of height $k$ if $i_0=k$.
\end{proposition}

\formal{\lean{HPC.HO16.ho16_prop19}, whose statement is the proposition \lean{HPC.HO16.HOProp19}; for the full (possibly disconnected) fold,
\lean{HPC.HO16.ho16_prop19_full}.}

\begin{proof}[Proof outline]
\emph{The fold.} Put $G_1=\pi(P_1)$, $G_2=\pi(P_2)$, $G_3=\side_{\pi(\gamma)}(\pi(S_{i_0}))$. By Proposition~\ref{prop:sides-char}(d) and (S3$'$), $\pi(P_1)=\side_{\pi(\mathcal E(P_1))}(\pi(P_2))$; by
Proposition~\ref{prop:ho10}, $G_3=\side_{\pi(\gamma)}(\pi(P_2))$. So Proposition~\ref{prop:ho11} with $A=\pi(P_2)$, $B_1=\pi(\mathcal E(P_1))$, $B_2=\pi(\gamma)$ shows that $(G_1,G_2,G_3)$ is a fold triple,
and
\begin{equation}\label{eq:19-frontiers}
\rfr G_1=\pi(\mathcal E(P_1)),\qquad \rfr G_3=\pi(\gamma),\qquad \rfr G_2=\pi(\mathcal E(P_1))\cup\pi(\gamma).
\end{equation}
Theorem~\ref{thm:fold-exists} gives a simple fold $F$.

\emph{The sets.} For $i\notin\set{i_0,i_0+1}$: $S_i'=\varphi^{\star-1}(S_i)$, $\alpha_i'=\varphi^{\star-1}(\alpha_i)$. At level $i_0$:
\[
S_{i_0}'=[\varphi^{\star-1}(P_1)\cap F_1^\star]\cup[\varphi^{\star-1}(P_2)\cap F_2^\star]\cup[\varphi^{\star-1}(S_{i_0})\cap F_3^\star],\qquad \alpha_{i_0}'=\varphi^{\star-1}(\alpha_{i_0}),
\]
and $\alpha_{i_0+1}'=\varphi^{\star-1}(\alpha_{i_0+1})\cap F_3^\star$. If $i_0<k$: $S_{i_0+1}'=\varphi^{\star-1}(S_{i_0+1})$, $\gamma'=\varphi^{\star-1}(\alpha_{i_0+1})\cap(F_1\cup F_2)^\star$,
$P_1'=\varphi^{\star-1}(P_2)\cap(F_1\cup F_2)^\star$ and $P_2'=\varphi^{\star-1}(S_{i_0})\cap(F_1\cup F_2)^\star$. In words: the new sheet $i_0$ runs along $P_1$ over $F_1$, back along $P_2$ over $F_2$, and
on along $S_{i_0}$ over $F_3$: the pit has been unwound. The part of sheet $i_0+1$ over $F_1\cup F_2$ becomes the new pit.

\emph{Checks.}
\begin{itemize}
\item For $i\neq i_0$, $\rfr\pi(S_i)=\pi(\alpha_i)\cup\pi(\beta_i)$ misses $\rfr G_2$ by (S5$'$), so Lemma~\ref{lem:ho18}(1) makes $S_i'$ straight with $\mathcal E(S_i')=\varphi^{\star-1}(\mathcal E(S_i))$.
\item At level $i_0$: over $\rfr F_1$ the pieces from $P_1$ and $P_2$ coincide, because $\mathcal E(P_1)\subseteq\mathcal E(P_2)\subseteq P_2$ and $\pi$ is injective on $P_2$; likewise over $\rfr F_3$
with $\gamma\subseteq\mathcal E(P_2)\cap\mathcal E(S_{i_0})$. So $S_{i_0}'$ is straight; $F_1\cup F_2$ is interior to $\pi(S_{i_0}')$, and the end points of $S_{i_0}'$ lie over the interior of $F_3$, as
$\pi(S_{i_0})$ agrees with $G_3$ near $\pi(\gamma)$ (Proposition~\ref{prop:sides-char}(c)).
\item $P_1'$ and $P_2'$: Lemma~\ref{lem:ho18}(2) with $S=P_2$ and $S=S_{i_0}$, and $\mathcal E(P_1')=\varphi^{\star-1}(\gamma)$, $\mathcal E(P_2')=\mathcal E(P_1')\cup\gamma'$, using (S6$'$).
\item Finiteness in (S2): a point of $G$ has at most three preimages in $F$. The neighbourhoods of (S3) pull back, through local homeomorphisms off $\rfr F_2$ and, at a point
$z\in\rfr F_3$, through the two spokes of $F$ at $z$ (one in $F_1\cup F_2$, one in $F_3$), which exist since $\varphi(z)$ is ordinary.
\item (S4): a component $C$ of $F\setminus\pi(S_i')$ maps into a component $\tilde C$ of $G\setminus\pi(S_i)$, with $\varphi(\rfr C)\subseteq\rfr\tilde C$, and consistency for $S_i$ applies; at
level $i_0+1$ one uses Proposition~\ref{prop:ho10} (second part) and the separation of Lemma~\ref{lem:ho6}(6).
\item (S5): if every point of $\rfr G_2$ is ordinary in $G$, a point $z\in F$ is ordinary in $F$ iff $\varphi(z)$ is ordinary in $G$; and preimages of disjoint sets are disjoint.
\item (S6$'$): $\alpha_{i_0+1}'$ lies over the interior of $F_3$.
\end{itemize}

\emph{Connectedness.} By Proposition~\ref{prop:ho7} there is a component $C$ of $F$ whose image meets $\rfr G_1$ and $\rfr G_3$; then $\varphi(C)=G$ and the pieces $F_i\cap C$ form a
connected simple fold. Every set of the new structure meets $C^\star$: for $i\neq i_0$ because $\varphi(C)=G$; $S_{i_0}'$, $P_1'$, $P_2'$ because each contains the points over
$\gamma$ in $(\rfr F_3)^\star$, and $C$ meets $\rfr F_3$ by Lemma~\ref{lem:ho6}(6) ($C$ meets both $F_1\setminus F_2$ and $F_3\setminus\rfr F_3$, since $G\setminus G_3$ and $G\setminus G_1$ are nonempty).
Restrict the structure to $C$ (Proposition~\ref{prop:stair-remarks}(a) and its broken analogue).
\end{proof}

\formal{The checks are spread over \lean{HPC/HO16/Prop19A.lean} to \lean{HPC/HO16/Prop19D.lean}; the fold is \lean{HPC.HO16.ho16_prop19_fold}; the sheets are
\lean{HPC.HO16.p19T}, with a companion definition for the end sets in \lean{HPC/HO16/P19Defs.lean}; ordinary points of folds: \lean{HPC.HO16.IsSimpleFoldOf.isOrdinaryPt_iff}.}

\begin{remark}[A false sentence on p.~208]\label{rem:p208}
In the (S4$'$) check for $P_2'$, Hoehn and Oversteegen take a component $D\subseteq F_1\cup F_2$ of $F\setminus\pi(P_2')$, its image in a component $\tilde D$ of $G\setminus\pi(S_{i_0})$, and assert
that $\rfr\tilde D\cap\pi(\alpha_{i_0})=\emptyset$ by (S6$'$). This is false in a legitimate configuration: take $G=[0,10]$, $k=3$, $i_0=2$, $5<a<6$,
\begin{itemize}
\item $S_1=[4,10]\times\set{0.1}$, $S_2=[(4,0.1),(6,0.5)]$, $S_3=[0,10]\times\set{0.9}$; $\alpha_2=\set{(4,0.1)}$, $\gamma=\set{(6,0.5)}$, all other $\alpha_i$ empty;
\item $P_2=[1,2]\times\set{0.3}\cup[(a,0.3),(6,0.5)]$, and $P_1$ a path over $[1,2]$ from $(1,0.3)$ to $(2,0.3)$ through $(1.5,0.2)$, together with $[(a,0.3),(10,0.2)]$.
\end{itemize}
Then $x=1.5\in\pi(P_2)\setminus\pi(S_2)$ lies in the component $\tilde D=[0,4)$ of $G\setminus\pi(S_2)$, and $\rfr\tilde D=\set4=\pi(\alpha_2)$. In the fold, the copies of $[1,2]$ in $F_1$ and $F_2$ form a
circle component $D$ of $F$ over $\tilde D$. The conclusion of the step survives, by a case split: if $\rfr D=\emptyset$ there is nothing to show; otherwise a point $z\in\rfr D$ has
$\varphi(z)\in\rfr\tilde D\cap\pi(P_1\cup P_2)$, which misses $\pi(\alpha_{i_0})$ by (S6$'$), so consistency puts $\rfr\tilde D$ inside $\pi(\beta_{i_0})$ or inside $\pi(\gamma)$.
\end{remark}

\formal{\lean{HPC.HO16.ho16_p208_false}.}

On p.~205 Hoehn and Oversteegen also print ``$P_1'\cap(\rfr F_3)^\star=\varphi^{\star-1}(\gamma_{i_0})$''; the right-hand side also contains the points over the $F_1$-copies of $\pi(\gamma)$, outside
$(\rfr F_3)^\star$, so only ``$\subseteq$'' holds, which is all that is used.

\section{Theorem 20 from Theorem 15}

\begin{theorem}[Hoehn--Oversteegen, Theorem~15]\label{thm:ho15}
Let $G$ be a connected PL graph, $M\subseteq G\times(0,1)$ a set separating $G\times\set0$ from $G\times\set1$ in $G^\star$, and $U\supseteq M$ open. Then some $S\subseteq U\cap G\times(0,1)$ has a
stairwell structure of odd height.
\end{theorem}

\formal{The statement is the proposition \lean{HPC.HO16.HOThm15}; it is proved as \lean{HPC.HO16.ho16_theorem15} (Chapter~\ref{ch:ladders}).}

Hoehn and Oversteegen state Theorem~15 for any graph. Some nonemptiness is needed: for $G=\emptyset$ and $M=\emptyset$ the separation holds with empty open sets, and no stairwell
lives in $\emptyset$. Connectedness is all that the proof of Theorem~20 uses.

\begin{theorem}\label{thm:lift-stairwell}
Let $G$ be a connected PL graph, $S\subseteq G^\star$ with a stairwell structure of odd height $k$, $X$ a hereditarily indecomposable compact metric space and $f\colon X\to G$ continuous.
For every $\varepsilon>0$ there is a continuous $h\colon X\to S$ with $\nrm{f(x)-\pi h(x)}<\varepsilon$.
\end{theorem}

\begin{proof}
Prove simultaneously the same statement for broken structures of odd height (pit at any level), by strong induction on the measure $(k+1)^2$ for a stairwell of height $k$ and
$(k+1)^2+k+1-i_0$ for a broken structure with a pit at $i_0$, quantified over the normed space $E$ (the induction passes from $E$ to $E\times\R$).
\begin{itemize}
\item Height $1$: $\pi|_S$ is a homeomorphism onto $G$ (Proposition~\ref{prop:stair-remarks}(b)); take $h=(\pi|_S)^{-1}\circ f$.
\item Height $k\ge3$: Proposition~\ref{prop:ho17} gives a broken structure of height $k-2$ with a pit at level~$1$, of smaller measure.
\item Broken, with $\gamma=\emptyset$ or $\mathcal E(P_1)=\emptyset$: Proposition~\ref{prop:rem17} gives a stairwell structure of height $k$ or $1$ on a subset of $S$, of smaller measure.
\item Broken, otherwise: Proposition~\ref{prop:ho19} gives a connected fold $F$ and $S'\subseteq\varphi^{\star-1}(S)$ with a structure of smaller measure in $F^\star$. Theorem~\ref{thm:ho8} gives
$g\colon X\to F$ with $\nrm{f(x)-\varphi(g(x))}<\varepsilon/2$; the induction hypothesis for $F$, $S'$, $g$ and $\varepsilon/2$ gives $h'\colon X\to S'$; put $h=\varphi^\star\circ h'$, which maps into $S$.
Since $\varphi$ is a projection, $\nrm{f(x)-\pi h(x)}\le\nrm{f(x)-\varphi g(x)}+\nrm{g(x)-\pi h'(x)}<\varepsilon$.
\end{itemize}
Oddness is preserved at each step, which excludes height $2$, where Proposition~\ref{prop:ho17} does not apply.
\end{proof}

\formal{\lean{HPC.HO16.ho16_lift_stairwell}, with the induction in \lean{HPC.HO16.ho16_lift_aux}.}

Hoehn and Oversteegen compose the folds explicitly, $h=(\varphi_1)^\star\circ\dots\circ(\varphi_n)^\star\circ\theta\circ g_n$; the recursion does the same without indexing them.

\begin{theorem}\label{thm:20-from-15}
Theorem~\ref{thm:ho15} implies Theorem~\ref{thm:ho20}.
\end{theorem}

\begin{proof}
$G$ has finitely many components $C$, each clopen; $M\cap C^\star$ separates $C\times\set0$ from $C\times\set1$ in $C^\star$ (with the same open sets), and $f^{-1}(C)$ is a clopen, hereditarily
indecomposable compactum. On each piece, Theorem~\ref{thm:ho15} gives $S\subseteq U\cap C\times(0,1)$ with a stairwell structure of odd height and Theorem~\ref{thm:lift-stairwell} gives $h$; glue
the pieces.
\end{proof}

\formal{\lean{HPC.HO16.ho16_theorem20_of_thm15}; with Theorem~15 as the only hypothesis, the classification step is
\lean{HPC.HO16.homeomorph_pseudoArc_of_indec_homogeneous_of_thm15}.}

\section*{Exercises}

\begin{exercise}
Let $G=\I$, $S_1=[0,\frac23]\times\set{\frac14}$, $S_2$ the segment from $(\frac23,\frac14)$ to $(\frac13,\frac12)$ and $S_3$ the segment from $(\frac13,\frac12)$ to $(1,\frac34)$.
Show that $\langle S_1,S_2,S_3\rangle$, with $\alpha_2=\set{(\frac23,\frac14)}$ and $\alpha_3=\set{(\frac13,\frac12)}$, is a stairwell structure of height $3$, and check Claim~14.1.
\end{exercise}

\begin{exercise}
For the stairwell of the previous exercise, carry out Proposition~\ref{prop:ho17} and one application of Proposition~\ref{prop:ho19}. Draw the fold and the new structure.
\end{exercise}

\begin{exercise}
Show that for a stairwell structure of even height on a connected $G$, $\pi(S)$ need not be $G$.
\end{exercise}

\begin{exercise}
Check the five conditions (S1$'$)--(S6$'$) for the configuration of Remark~\ref{rem:p208}.
\end{exercise}

\begin{exercise}
Why must the measure in the proof of Theorem~\ref{thm:lift-stairwell} depend on $i_0$?
\end{exercise}
